\chapter{Hydrocele Repair}

\section*{Introduction \& Clinical Indications}
Hydrocele repair, also known as hydrocelectomy, is a surgical procedure performed to remove or obliterate a hydrocele, which is an abnormal collection of serous fluid within the tunica vaginalis, the sac surrounding the testicle. This condition typically manifests as a painless swelling in the scrotum, which can vary significantly in size and consistency, often causing discomfort or cosmetic concerns. The primary goal of the surgery is to eliminate scrotal swelling, alleviate associated symptoms such as heaviness or pain, and prevent potential complications like infection or interference with daily activities. While often benign, surgical intervention is indicated when the hydrocele becomes symptomatic, progressively enlarges, or when there is diagnostic uncertainty regarding an underlying testicular pathology, such as a tumor. The procedure aims for a definitive resolution of the fluid accumulation, restoring normal scrotal anatomy and function.

\begin{itemize}
  \item Hydrocelectomy
  \item Excision of hydrocele sac
  \item Eversion of hydrocele sac
  \item Plication of hydrocele sac
  \item Lord's plication procedure
  \item Jaboulay's procedure
  \item Winkelmann's procedure
\end{itemize}

A hydrocele results from an imbalance between the secretion and absorption of fluid by the mesothelial lining of the tunica vaginalis. In adults, this is typically a non-communicating hydrocele, where the processus vaginalis has obliterated, and fluid accumulates due to inflammation, trauma, infection, or idiopathic causes. In children, a communicating hydrocele occurs due to a patent processus vaginalis, a remnant of the embryonic peritoneal outpouching, allowing peritoneal fluid to enter the scrotum. The surgical principle involves either excising the hydrocele sac entirely or everting/plicating it to permanently prevent further fluid accumulation by altering the secretory surface or reducing the sac volume.

The rationale behind surgical intervention is to definitively address the underlying mechanism of fluid accumulation. Excision removes the fluid-producing lining, eversion exposes the secretory surface to the absorptive scrotal tissues, and plication reduces the surface area available for fluid secretion and promotes absorption. The technical goals are to achieve complete drainage of the fluid, permanently prevent its reaccumulation by altering the tunica vaginalis, and meticulously preserve the integrity of the spermatic cord structures, including the testicular artery, veins, and vas deferens, to ensure testicular viability and function. This approach aims to alleviate symptoms, improve cosmesis, and prevent potential complications associated with persistent scrotal swelling.

The treatment of hydroceles dates back to ancient times, with early methods involving aspiration and injection of sclerosing agents, often with high recurrence rates and significant complications such as infection and chemical peritonitis. Significant advancements in surgical technique emerged in the 19th century, moving towards open surgical repair. In 1882, Dr. Carl Winkelmann described the eversion technique, where the hydrocele sac was opened, everted, and sutured behind the testis, exposing its inner lining to the scrotal tissues to promote absorption. Shortly after, in 1893, Dr. Louis Jaboulay introduced the excision technique, advocating for the complete removal of the hydrocele sac, leaving only a small cuff around the epididymis and testis. These open surgical approaches became the standard of care and remain foundational to modern hydrocelectomy.

In the mid-20th century, Dr. Peter Lord introduced the plication technique (Lord's procedure), which gained popularity for its simplicity and reduced dissection, particularly in cases of thin-walled hydroceles. This technique involves opening the sac and pleating or folding its walls upon themselves with sutures, effectively reducing its volume without extensive dissection around the epididymis and vas deferens, potentially reducing postoperative complications. The evolution of anesthesia, sterile techniques, and refined surgical instruments further improved outcomes, making hydrocele repair a safe and highly effective procedure with low morbidity and recurrence rates.

Hydrocele repair is indicated for a variety of clinical scenarios, primarily when the hydrocele becomes symptomatic or poses diagnostic challenges. Specific indications include:

\begin{itemize}
  \item Significant scrotal swelling causing discomfort, pain, or a heavy sensation that interferes with daily activities.
  \item Large hydroceles that impede ambulation, sexual function, or the wearing of clothing.
  \item Cosmetic concerns due to the size and appearance of the scrotum, impacting quality of life.
  \item Diagnostic uncertainty, particularly when a hydrocele obscures palpation of the testis, necessitating surgical exploration to rule out an underlying testicular tumor or other pathology.
  \item Recurrent hydroceles after previous aspiration attempts, which typically have high recurrence rates.
  \item Secondary hydroceles associated with epididymitis, orchitis, trauma, or testicular torsion that persist after resolution of the primary condition.
  \item Communicating hydroceles in children that persist beyond 12-18 months of age, are symptomatic, or are associated with an inguinal hernia.
  \item Rapidly enlarging hydroceles or those causing pressure symptoms on the testis.
\end{itemize}

Hydrocele repair is generally considered the primary definitive treatment for symptomatic hydroceles that do not resolve spontaneously or are recurrent after less invasive measures. For adult non-communicating hydroceles, surgery is the first-line intervention once conservative management, such as observation for small, asymptomatic hydroceles, has failed or is deemed inappropriate due to size or symptoms. It is the preferred approach for long-term resolution.

In children with communicating hydroceles, surgery is also a primary treatment, typically performed if the hydrocele persists beyond infancy (usually 12-18 months of age) or is associated with a symptomatic inguinal hernia. Aspiration of hydroceles is considered a secondary, temporary measure, often used for diagnostic purposes or in patients who are poor surgical candidates due to significant comorbidities. However, aspiration carries a high recurrence rate and does not offer a definitive solution, often leading to repeat procedures or eventual surgical repair. Therefore, hydrocele repair serves as the primary surgical solution for long-term resolution of the condition, offering superior efficacy and lower recurrence compared to non-surgical alternatives.

\section*{Relevant Surgical Anatomy}
The scrotum is a musculocutaneous pouch containing the testes, epididymides, and the distal portions of the spermatic cords. Its wall consists of several layers: skin, dartos fascia (containing dartos muscle), external spermatic fascia, cremasteric fascia (containing cremaster muscle), and internal spermatic fascia. Deep to these layers lies the tunica vaginalis, a serous membrane derived from the peritoneum, which has two layers: a parietal layer lining the internal spermatic fascia and a visceral layer intimately covering the testis and epididymis. The space between these two layers, the tunica vaginalis cavity, normally contains a small amount of lubricating fluid. A hydrocele forms when there is an excessive accumulation of fluid within this cavity.

The testis is an ovoid organ responsible for spermatogenesis and hormone production, located within the tunica vaginalis. Posteriorly, the epididymis, a coiled tube, is attached to the testis and serves as a site for sperm maturation and storage, eventually leading to the vas deferens. The spermatic cord, which enters the scrotum through the inguinal canal, contains vital structures including the testicular artery (supplying the testis), the pampiniform plexus of veins (draining the testis and epididymis), the vas deferens (transporting sperm), lymphatic vessels, and nerves (genitofemoral nerve, ilioinguinal nerve). During hydrocele repair, meticulous dissection is required to separate the hydrocele sac from these delicate structures, particularly the spermatic cord, to prevent injury and preserve testicular viability and function.

\section*{Preoperative Workup \& Preparation}
Preoperative preparation for hydrocele repair involves a comprehensive assessment to ensure patient safety, confirm the diagnosis, and optimize surgical outcomes. A thorough medical history and physical examination are performed to identify any comorbidities that might increase surgical risk. Routine blood tests, including a complete blood count (CBC) to assess for anemia or infection, a coagulation profile (PT/INR, PTT) to evaluate bleeding risk, and a basic metabolic panel to check kidney function and electrolytes, are usually ordered.

Key preparatory steps include:
\begin{itemize}
  \item \textbf{Diagnostic Confirmation:} A scrotal ultrasound is essential to confirm the diagnosis of a hydrocele, assess its size and characteristics, and crucially, to rule out any underlying testicular pathology such as a tumor, epididymitis, or hernia, which would significantly alter the surgical plan.
  \item \textbf{Anesthesia Clearance:} Patients undergo an anesthesia evaluation to assess their fitness for general or regional anesthesia. This may involve an electrocardiogram (ECG) for cardiac assessment and a chest X-ray for older patients or those with significant pulmonary comorbidities.
  \item \textbf{Medication Review and Adjustment:} Patients are instructed to discontinue blood-thinning medications (e.g., aspirin, NSAIDs, warfarin, novel oral anticoagulants) for a specified period (typically 5-7 days) before surgery to minimize bleeding risk. Other medications are reviewed and adjusted as necessary.
  \item \textbf{NPO Status:} Patients must adhere to "nil per os" (NPO) guidelines, meaning no food or drink for at least 6-8 hours prior to surgery, to prevent aspiration during anesthesia induction.
  \item \textbf{Antibiotic Prophylaxis:} Prophylactic antibiotics, typically a first or second-generation cephalosporin, are administered intravenously within 60 minutes prior to skin incision, especially if there is a higher risk of infection or for longer procedures.
  \item \textbf{Informed Consent:} A detailed discussion with the surgeon regarding the procedure, its rationale, potential risks, expected benefits, and alternative treatment options is conducted, followed by obtaining the patient's informed consent.
\end{itemize}

\textbf{Absolute Contraindications:}
\begin{itemize}
  \item Active scrotal or systemic infection: Surgery should be postponed until the infection is fully treated and resolved to prevent wound complications and sepsis.
  \item Uncorrected coagulopathy or severe bleeding disorder: Significant risk of uncontrolled hemorrhage during or after surgery.
  \item Severe systemic illness or unstable medical condition: Patients who are not medically fit for anesthesia and surgery due to conditions like uncontrolled heart failure, recent myocardial infarction, or severe respiratory compromise.
  \item Hydrocele secondary to testicular malignancy: In such cases, an orchiectomy (testicle removal) is typically indicated, not a simple hydrocele repair, to ensure complete oncological resection.
\end{itemize}

\textbf{Relative Precautions:}
\begin{itemize}
  \item Small, asymptomatic hydroceles: Observation may be appropriate, especially if the hydrocele is not causing discomfort, functional impairment, or cosmetic concern.
  \item Hydroceles in infants that may resolve spontaneously: Communicating hydroceles in children often resolve by 12-18 months of age; surgery is typically deferred until then unless symptomatic or associated with an inguinal hernia.
  \item Patients with significant comorbidities: Careful risk-benefit assessment is necessary. Less invasive options (e.g., aspiration) might be considered if surgery poses excessive risk, though these are not definitive treatments.
  \item Previous scrotal surgery: May increase the technical difficulty of dissection due to scar tissue and adhesions, requiring a more cautious approach.
  \item Very large or long-standing hydroceles: These can be associated with thinning of the scrotal skin and potential testicular atrophy, requiring careful handling and sometimes more extensive skin resection.
\end{itemize}

\section*{Surgical Technique \& Procedure}
Hydrocele repair is typically performed under general or regional (spinal) anesthesia, often in an outpatient setting. The patient is positioned supine, and the scrotal and inguinal areas are prepped with an antiseptic solution and draped in a sterile fashion.

The surgical approach depends on the type and size of the hydrocele, as well as patient age:
\begin{itemize}
  \item \textbf{Scrotal Approach:} For most adult non-communicating hydroceles, a transverse scrotal incision is made directly over the hydrocele. This allows direct access to the tunica vaginalis.
  \item \textbf{Inguinal Approach:} For communicating hydroceles in children or when an associated inguinal hernia is suspected, an inguinal incision is preferred. This allows for ligation of the patent processus vaginalis at the internal ring, similar to an inguinal hernia repair, in addition to managing the scrotal component.
\end{itemize}

Key operative steps typically include:
\begin{enumerate}
  \item \textbf{Incision and Exposure:} A transverse scrotal incision is made through the skin and dartos muscle. The underlying layers are carefully dissected to expose the tunica vaginalis, which appears as a bluish, tense sac.
  \item \textbf{Hydrocele Sac Isolation:} The hydrocele sac is carefully identified and meticulously isolated from the surrounding scrotal tissues and the delicate structures of the spermatic cord (vas deferens, testicular vessels). This requires careful blunt and sharp dissection.
  \item \textbf{Sac Incision and Fluid Aspiration:} The hydrocele sac is incised, and the serous fluid is aspirated. The fluid is typically clear or straw-colored and may be sent for pathological analysis if there is any suspicion of malignancy.
  \item \textbf{Definitive Repair Technique:} One of the following techniques is performed:
    \begin{itemize}
      \item \textbf{Excision (Jaboulay's or Bottle Technique):} The majority of the hydrocele sac is excised, leaving only a small cuff of tunica vaginalis (approximately 1 cm) around the epididymis and testis. The edges of the remaining cuff are then oversewn with absorbable sutures to achieve meticulous hemostasis and prevent re-eversion.
      \item \textbf{Eversion (Winkelmann's Technique):} The opened hydrocele sac is everted and sutured behind the testis using absorbable sutures. This exposes the inner lining of the sac to the absorptive scrotal tissues, promoting fluid reabsorption and preventing reaccumulation.
      \item \textbf{Plication (Lord's Technique):} The sac is opened, and its walls are pleated or folded upon themselves with non-absorbable sutures, effectively reducing its volume and surface area without extensive dissection or excision. This technique is often preferred for thin-walled hydroceles.
    \end{itemize}
  \item \textbf{Hemostasis and Testicular Repositioning:} Meticulous hemostasis is achieved throughout the procedure to minimize postoperative bleeding and hematoma formation. The spermatic cord structures and testis are carefully inspected for integrity and returned to their anatomical position within the scrotum.
  \item \textbf{Wound Closure:} The surgical wound is closed in layers using absorbable sutures. The dartos muscle is approximated, followed by skin closure. A small Penrose drain may occasionally be placed in large hydroceles or if significant oozing is anticipated, though it is not routinely used. A scrotal support dressing is applied.
\end{enumerate}

\section*{Complications \& Risks}
While hydrocele repair is generally a safe and effective procedure, like any surgical intervention, it carries potential risks and complications:

\begin{itemize}
  \item \textbf{General Surgical Risks:}
    \begin{itemize}
      \item \textit{Anesthesia-related risks:} These include adverse reactions to medications, respiratory depression, cardiac events, and postoperative nausea and vomiting.
      \item \textit{Bleeding and Hematoma:} Accumulation of blood within the scrotum, which can cause significant swelling, pain, and may occasionally require further surgical intervention for evacuation.
      \item \textit{Infection:} Wound infection at the incision site, or rarely, more severe infections such as epididymo-orchitis or Fournier's gangrene.
      \item \textit{Pain:} Postoperative pain, typically managed with oral analgesics, but can occasionally be chronic.
    \end{itemize}
  \item \textbf{Procedure-Specific Risks:}
    \begin{itemize}
      \item \textit{Scrotal Edema and Swelling:} Common postoperatively, often resolving over several weeks to months, but can be persistent.
      \item \textit{Recurrence of Hydrocele:} Although uncommon (2-5\%), the hydrocele can reaccumulate, potentially requiring repeat surgery. This risk can vary with the technique used and the underlying cause.
      \item \textit{Injury to Testis or Spermatic Cord Structures:} Rare but serious complications include damage to the testicular artery, veins (pampiniform plexus), or vas deferens, potentially leading to testicular atrophy, infertility, or chronic pain.
      \item \textit{Epididymitis/Orchitis:} Inflammation or infection of the epididymis or testis, which can cause pain and swelling.
      \item \textit{Nerve Injury:} Damage to scrotal nerves can result in areas of numbness, altered sensation, or chronic neuropathic pain.
      \item \textit{Testicular Atrophy:} Very rare, but possible if the blood supply to the testis is significantly compromised during dissection.
      \item \textit{Chronic Scrotal Pain:} Persistent pain in the scrotum, which can be debilitating and difficult to treat.
      \item \textit{Cosmetic Dissatisfaction:} While the goal is to improve cosmesis, some patients may have concerns about scarring or residual scrotal asymmetry.
    \end{itemize}
\end{itemize}

\section*{Postoperative Care \& Recovery}
Immediately after hydrocele repair, the patient is transferred to the Post-Anesthesia Care Unit (PACU) for close monitoring as they recover from anesthesia. Pain management is initiated, and patients are typically discharged home the same day or the following morning, depending on their recovery, the type of anesthesia used, and the complexity of the procedure. Patients are advised to arrange for a responsible adult to drive them home.

During the first 24-48 hours, patients are advised to rest, apply ice packs to the scrotum intermittently to reduce swelling and discomfort, and wear a scrotal support (jockstrap or tight-fitting underwear) continuously to provide comfort and minimize edema. Mild to moderate pain is expected and effectively managed with prescribed oral analgesics. Scrotal swelling and bruising are common and gradually subside over several weeks. For the first 1-2 weeks, patients should avoid strenuous activities, heavy lifting (typically over 10-15 pounds), and sexual activity. Wound care involves keeping the incision site clean and dry; patients are usually advised to avoid baths, hot tubs, or swimming until the incision is fully healed, typically 10-14 days. Most individuals can return to light work or school within 1 week, with a gradual return to full activity, including exercise and sexual activity, over 4-6 weeks. Long-term recovery involves continued monitoring for any signs of recurrence or complications, though these are rare.

During hydrocele repair, the surgeon typically observes a tense, fluid-filled sac (the tunica vaginalis) surrounding the testis. The fluid is usually clear, straw-colored, or occasionally amber. The thickness and characteristics of the hydrocele sac walls are noted, as this can influence the choice of surgical technique (e.g., thin walls for plication, thicker walls for excision). The testis and epididymis are inspected for any abnormalities once the fluid is drained. Any adhesions between the sac and surrounding structures or the testis are carefully lysed.

The pathology report on the excised hydrocele sac (if excision is performed) typically confirms the presence of serous fluid and benign mesothelial lining, consistent with a hydrocele. The tissue examination rules out any unexpected malignancy or other pathological processes within the sac wall. In cases of communicating hydroceles, the patent processus vaginalis is identified and ligated. The discharge summary will detail the specific procedure performed, any intraoperative findings (e.g., fluid volume, sac characteristics, condition of the testis), and postoperative instructions. The absence of any malignant cells or unusual findings on pathology provides reassurance regarding the benign nature of the condition.

A normal, successful postoperative course following hydrocele repair is characterized by a progressive reduction in scrotal swelling and discomfort, leading to a return to normal scrotal contour and function. Initial pain, swelling, and bruising are expected but should gradually improve over the first few days to weeks. Pain is typically well-controlled with oral analgesics and usually resolves within 1-2 weeks. Scrotal edema, which can be significant initially, will slowly subside over 4-6 weeks, though some residual firmness or minor swelling may persist for a few months. The incision site should heal cleanly, with sutures (if non-absorbable) removed at the first follow-up visit. Patients should experience a gradual return to normal daily activities, with full physical activity typically resumed by 4-6 weeks post-surgery. The primary indicator of a successful outcome is the complete resolution of the hydrocele and the alleviation of preoperative symptoms such as discomfort, heaviness, or cosmetic concerns, without significant complications or recurrence.

Postoperative care and follow-up are crucial to monitor recovery, ensure proper wound healing, and detect any potential complications or recurrence.

\begin{itemize}
  \item \textbf{Postoperative Clinic Visit:} A follow-up appointment with the surgeon is typically scheduled 1-2 weeks after surgery. During this visit, the wound healing is assessed, any non-absorbable sutures or staples (if used) are removed, and the scrotum is examined for early complications like infection, hematoma, or excessive swelling.
  \item \textbf{Activity Resumption Guidance:} Patients receive detailed guidance on gradually resuming physical activities, including exercise, sports, and sexual activity, usually over a period of 4-6 weeks. Restrictions on heavy lifting are emphasized.
  \item \textbf{Wound Care Instructions:} Patients are instructed on how to care for the incision site, typically involving keeping it clean and dry, and avoiding immersion in water until fully healed.
  \item \textbf{Warning Signs Education:} Patients are educated on warning signs that necessitate prompt contact with their doctor, such as increasing pain unresponsive to medication, fever (above 101.5$^{\circ}$F or 38.6$^{\circ}$C), redness, warmth, pus drainage from the incision, or sudden re-swelling of the scrotum.
  \item \textbf{Long-term Monitoring:} While recurrence is low, patients are advised to perform regular self-examinations and report any new scrotal swelling or discomfort. No routine imaging or laboratory tests are typically required unless complications or recurrence are suspected.
\end{itemize}
Adherence to these instructions is vital for optimal healing and long-term success.

\section*{Outcomes \& Prognosis}
Hydroceles are a common urological condition with varying incidence rates depending on age and type. Communicating hydroceles are frequently observed in neonates, affecting approximately 1-5\% of full-term male infants, often resolving spontaneously by 12-18 months of age as the processus vaginalis obliterates. Adult non-communicating hydroceles are more prevalent in men over 40, with an estimated incidence of 1\% in the general adult male population. The exact etiology of adult hydroceles is often idiopathic, but they can be secondary to trauma, infection (e.g., epididymitis, orchitis), testicular torsion, or tumors. The frequency of surgical repair is directly related to the prevalence of symptomatic cases that require intervention. Mortality associated with hydrocele repair is exceedingly rare, primarily linked to severe anesthesia complications in high-risk patients with significant comorbidities. Morbidity is generally low, with minor complications such as scrotal swelling, bruising, and pain being the most common, while serious complications like infection, significant hematoma, or recurrence occur in a small percentage of cases (typically 2-5\%).

The outcomes of hydrocele repair are generally excellent, with high success rates and a very favorable prognosis. Surgical correction effectively resolves the scrotal swelling and associated symptoms in 90-95\% of cases, providing lasting relief. Functional outcomes are typically very good, with patients experiencing significant improvement in comfort, resolution of heaviness, improved cosmetic appearance, and the ability to resume normal activities without limitation. Recurrence rates are low, generally ranging from 2-5\%, and are often associated with the specific surgical technique used (e.g., higher recurrence with aspiration, lower with excision/eversion), the size of the original hydrocele, or underlying contributing factors that were not fully addressed.

Prognostic factors for a successful outcome include meticulous surgical technique, complete excision or effective eversion/plication of the sac, appropriate patient selection, and adherence to postoperative care instructions. While rare, complications such as chronic pain or testicular atrophy can negatively impact prognosis, but the vast majority of patients achieve a complete and lasting resolution of their hydrocele with minimal long-term issues. The overall quality of life for patients undergoing successful hydrocele repair is significantly improved.

\section*{References}
\begin{enumerate}
  \item MedlinePlus. Hydrocele repair. U.S. National Library of Medicine; 2024. Available from: \url{https://medlineplus.gov/ency/article/002970.htm}
  \item Wein AJ, Kavoussi LR, Partin AW, Peters CA, editors. Campbell-Walsh-Wein Urology. 12th ed. Philadelphia, PA: Elsevier; 2021.
  \item American Urological Association. AUA Guideline on the Management of Scrotal Masses. 2023. Available from: \url{https://www.auanet.org/guidelines-and-quality/guidelines/scrotal-masses-guideline}
  \item Smith's General Urology. 19th ed. New York, NY: McGraw-Hill Education; 2020.
\end{enumerate}

\clearpage
